% Standalone wrapper to allow compiling this chapter directly.
\documentclass[12pt]{report}
\usepackage[utf8]{inputenc}
\usepackage[margin=1in]{geometry}
\usepackage{graphicx}
\usepackage{hyperref}
\usepackage{natbib}

\begin{document}

\chapter{Growth and Evolution of Windows OS Over the Past Three Decades}
\label{chap1:ch:windows_evolution}

\section{Introduction}
\label{chap1:sec:introduction}
As stated by Microsoft, Windows 1.0 was the very first version of Windows OS ever released[cite: 3]. Co-Founders of Microsoft, Bill Gates and Paul Allen had begun the development of the operating system in 1981, with official release taking place in November 1985[cite: 4]. The GUI revolution motivated Microsoft to develop an OS which had a visual interface which could compete with other emerging tech companies such as Apple inc and VisiCorp[cite: 4].

\begin{figure}[htbp]
    \centering
    % \includegraphics[width=0.8\textwidth]{images/chap1_windows_1_0.png}
    \caption{Placeholder: Visual interface of Windows 1.0.}
    \label{chap1:fig:windows_1_0}
\end{figure}

\section{Foundational Era: Windows 95 to Windows XP (1995 – 2005)}
\label{chap1:sec:foundational_era}
According to The Guardian, the transition from the more complex MS-DOS command line to the more consumer friendly Graphic User Interface (GUI) took place during this era, with this transition first being seen with the release of Windows 95[cite: 6]. The need for users to memorize text commands was eliminated as the visual workspace was introduced[cite: 6]. Internet access was integrated with Internet Explorer on Windows 98, along with the FAT32 file system which improved storage efficiency and reliability on the computer system[cite: 7].

As noted by ZDNET, in 2001, Windows XP was released[cite: 8]. Before Windows XP, Microsoft had separate ‘Home’ and ‘Business’ versions[cite: 8]. On Windows XP, both these versions were merged onto the stable NT kernel[cite: 9]. Windows integrated more features into the OS around this time such as the Luna interface which gave the system more visual and ‘bubble-like’ visuals, giving it a more modern feel at the time[cite: 10]. Wi-Fi support was also introduced as it was becoming mainstream at the time, and system recovery was improved to help users to recover data without having to reformat the system entirely[cite: 11].

\begin{figure}[htbp]
    \centering
    % \includegraphics[width=0.8\textwidth]{images/chap1_windows_xp.png}
    \caption{Placeholder: Windows XP Luna interface and start menu.}
    \label{chap1:fig:windows_xp}
\end{figure}

\section{Visual and Security Overhaul: Windows Vista and Windows 7 (2006 – 2011)}
\label{chap1:sec:visual_security}
Around this time, as highlighted by tech historians, Microsoft was working on implementing more high-end graphics into the Windows operating system[cite: 13]. Graphics such as the Aero Interface which introduced translucent window borders and hovering over the taskbar to see window previews were developed and implemented[cite: 14]. However, as the quality of the graphics improved, performance of the operating system began to decline[cite: 15]. User Account Control (UAC) is another important feature which was implemented around this time as a key security feature which asks for permission before making any system changes[cite: 16].

Windows 7 was then released and rectified the system issues which were presented by Windows Vista[cite: 17]. It was a more polished and enjoyable version of Windows Vista[cite: 18]. Snap was introduced to allow for easier window resizing and Libraries was also added, which made file management and organization an easier and more enjoyable experience for the user[cite: 19].

\begin{figure}[htbp]
    \centering
    % \includegraphics[width=0.8\textwidth]{images/chap1_windows_7.png}
    \caption{Placeholder: Windows 7 Aero interface and Snap feature.}
    \label{chap1:fig:windows_7}
\end{figure}

\section{The Mobile and Cloud Shift: Windows 8 and Windows 10 (2012 – 2020)}
\label{chap1:sec:mobile_cloud}
As touch-screen smartphones were rising in popularity, Microsoft needed to adapt[cite: 21]. With Windows 8, a full-screen tiled interface was designed and introduced[cite: 22]. This was considered a very risky and controversial move by Microsoft as it was a very big redesign of the interface[cite: 23]. One Drive was integrated which made cloud storage a core part of the operating system[cite: 24]. Windows 10 rectified the design “mistakes” of Windows 8 by reintroducing the original start menu while keeping some touch friendly features[cite: 25].

\begin{figure}[htbp]
    \centering
    % \includegraphics[width=0.8\textwidth]{images/chap1_windows_8.png}
    \caption{Placeholder: Windows 8 tiled full-screen interface.}
    \label{chap1:fig:windows_8}
\end{figure}

\section{The AI Era: Windows 11 and the Future (2021 – The Present)}
\label{chap1:sec:ai_era}
As stated by Indigo Software Company, current versions of Windows are focusing more on improving productivity, hybrid work and properly integrating artificial intelligence[cite: 27]. As AI is now at system level, AI is being used as a tool to improve user experience and increase the ease and efficiency of task completion by users[cite: 28]. Neural Processing Units (NPUs) are dedicated hardware which is installed in modern computers to run AI tasks such as real time translation and background noise removal[cite: 29].

\begin{figure}[htbp]
    \centering
    % \includegraphics[width=0.8\textwidth]{images/chap1_windows_ai.png}
    \caption{Placeholder: Integration of AI tools in modern Windows environments.}
    \label{chap1:fig:windows_ai}
\end{figure}

\section{Conclusion}
\label{chap1:sec:conclusion}
Implementation of AI has made the operating system more proactive rather than reactive as AI suggests files and layouts based on the task a user is completing instead of the user doing that themselves[cite: 31]. This reduces the workload for users and further increases task completion[cite: 32]. Microsoft also introduced technologies such as CompactOS, which uses LZX compression to shrink system files by up to 2 to 6 GB, allowing the OS to run more efficiently on devices with smaller SSDs[cite: 33].

\section*{References}
\label{chap1:sec:references}
\begin{itemize}
    \item Microsoft. 2026. Microsoft is Born - The Early History. [Online]. Redmond: Microsoft News. [cite: 35]
    \item The Guardian. 2014. From Windows 1 to 10: 29 Years of Evolution. [Online]. London: Guardian News \& Media. [cite: 36]
    \item Computer History Museum. 2026. The Mother of All Demos. [Online]. Mountain View: Computer History Museum. [cite: 37]
    \item Indigo Software. 2026. AI Features in Windows 2026 That Boost Productivity. [Online]. New York: Indigo Software Company. [cite: 38]
    \item ZDNET. 2026. 30 Years of Windows: 10 Milestones. [Online]. New York: Red Ventures. [cite: 39]
\end{itemize}

\end{document}
